% Generated from project outputs. Do not edit by hand.
\begin{table}[!htbp]
\centering
\begin{threeparttable}
\caption{Labor income by subsample: 2016-2024}
\label{tab:descriptive-2016-2024-subsamples-y1}
\small
\setlength{\tabcolsep}{4pt}
\renewcommand{\arraystretch}{1.15}
\begin{tabularx}{\linewidth}{@{}Xrrrrr@{}}
\toprule
Sample & N & Mean t & SD t & Mean t+1 & SD t+1 \\
\midrule
General & 552 & 2,712.6 & 1,677.9 & 2,907.8 & 1,692.0 \\
Men & 552 & 3,046.6 & 1,937.0 & 3,264.3 & 1,932.0 \\
Women & 552 & 2,162.5 & 1,498.6 & 2,353.7 & 1,577.9 \\
Urban & 552 & 2,962.4 & 1,721.1 & 3,141.0 & 1,716.7 \\
Rural & 547 & 2,208.2 & 1,556.0 & 2,436.6 & 1,695.7 \\
Indigenous & 547 & 2,191.5 & 1,587.2 & 2,465.0 & 1,666.5 \\
Non-Indigenous & 460 & 2,882.8 & 1,765.7 & 3,023.1 & 1,726.2 \\
Bottom 99 percent & 552 & 2,507.0 & 1,332.9 & 2,694.0 & 1,349.4 \\
Bottom 90 percent & 552 & 2,057.7 & 921.3 & 2,207.8 & 916.1 \\
Bottom 50 percent & 548 & 1,101.8 & 465.6 & 1,194.7 & 460.4 \\
\bottomrule
\end{tabularx}
\begin{tablenotes}[flushleft]
\footnotesize
\item Monthly income in quetzales. Unweighted descriptive statistics of survey-weighted cohort means. N counts cohort-transition rows, not individuals or independent cohorts. SD is dispersion across cohort means, not the standard error of a mean. Six annual transitions are pooled; the 2019-2021 gap is excluded. Both incomes must be observed. Subsamples overlap; bottom-income groups are cumulative.
\end{tablenotes}
\end{threeparttable}
\end{table}
